
\documentclass{ifcolog}

\usepackage{makeidx}


\title{The IFCoLog Class Style:\protect \\ Guide for Authors}

\titlerunning{IFCoLog Class Style}

%\addauthor{IFCoLog}{International Federation of Computational Logic}
\addauthor[michael.gabbay@kcl.ac.uk]
  {Michael Gabbay}
  {King's College London, Strand WC2R 2LS, UK}
\addauthor[m.gabbay@hw.ac.uk]
  {Murdoch Gabbay}
  {Heriot-Watt University Edinburgh, EH14 4AS, UK}

\authorrunning{IFCoLog}

\begin{document}
\maketitle

\begin{abstract}
A simple guide for using the IFColog class file. 
Version of February 2018.
\end{abstract}



\section{About the class file}
\label{about}


\begin{enumerate*}
\item
The default paper size is 7 by 10 inches. With a print area of 140 by 210 millimetres.
\item
The base font is Computer Modern, 10pt, and the {\sf sans-serif} font is {\sf Helvetica}.
\item
The reference list is condensed. The default bibliography style is \texttt{plain}.
\item
To insert images, pdfs etc, use the \texttt{graphicx} package.
For good printing quality images should be at least 300dpi (dots per inch).
\item
Note that page numbers will be \emph{changed} when papers are collected into a single volume.
\item
Figures are surrounded by a box.
\end{enumerate*}

\section{Preamble}
\label{preamble}

The front matter of an article is slightly different from others. Authors are added by a series of \verb+\addauthor+ commands, in addition to the \verb+\titlerunning+ and \linebreak \verb+\authorrunning+
commands, see Figure~\ref{fig1}.

\begin{figure}[t]
\begin{verbatim}
\title{The IFCoLog Class File: A Brief Guide}

\titlerunning{IFCoLog Class File}

\addauthor[michael.gabbay@kcl.ac.uk]
  {Michael Gabbay}
  {King's College London, Strand WC2R 2LS, UK}
\addauthor[m.gabbay@hw.ac.uk]
  {Murdoch Gabbay}
  {Heriot-Watt University Edinburgh, EH14 4AS, UK}

\authorrunning{$\backslash$authornames}

\begin{document}
\maketitle

\begin{abstract}
  A simple guide for using the IFColog class file.
\end{abstract}
\end{verbatim}
\caption{The preamble for this document}\label{fig1}
\end{figure}

The message reminding authors to thank the referees can be removed by placing the \verb+\titlethanks+ command in the preamble, followed by a message thanking the referees. 
To have no message use \verb+\titlethanks{}+.

\section{Section headings and capitalisation}
\label{section-headings}

Section and paragraph headings are invoked via the standard
commands, such as
\verb+\section+,
\verb+\subsection+,
\verb+\subsubsection+, and
Paper titles should follow the style commonly referred to as `Title Case' for headings:

\begin{enumerate*}
\item
Capitalise the first word in a title, regardless of part of speech.
\item
Capitalise all nouns, pronouns, verbs, adjectives, adverbs.
\item
Lowercase `to' as part of an infinitive.
\item
Lowercase all articles, prepositions and coordinating conjunctions.
\item
Other section headings should have just the first word capitalised.
\end{enumerate*}

\section{Overfull hboxes}

IFCoLog class file uses the \texttt{draft} option. A text overflow within any environment will be highlighted by a thick black line in the right margin. In the case of large tables or complex figures, use the \verb+\scalebox+ command to reduce sizes. For example in \texttt{quote} environment:
\begin{quote}
All work and no play makes Jack a dull boy. All work and no play makes Jack a dull boy. All work and no play makes Jack a dull boy. All work and no play makes Jack a dull boy.
All~work~and~no play~makes~Jack~a~dull~boy. All~work~and~no play~makes~Jack a~dull~boy. All~work~and~no play~makes~Jack a~dull~boy.
\end{quote}

\end{document}

